\documentclass[conference]{IEEEtran}
\IEEEoverridecommandlockouts
\usepackage{cite}
\usepackage{amsmath,amssymb,amsfonts}
\usepackage{algorithmic}
\usepackage{graphicx}
\usepackage{textcomp}
\usepackage{xcolor}
\usepackage[utf8]{inputenc} % Para reconocer tildes y caracteres en español
\usepackage[spanish,es-tabla]{babel} % Para los títulos en español y usar "Tabla" en vez de "Cuadro"
\usepackage{booktabs} % Fundamental para crear tablas estilo APA 7 (sin líneas verticales)

\def\BibTeX{{\rm B\kern-.05em{\sc i\kern-.025em b}\kern-.08em
    T\kern-.1667em\lower.7ex\hbox{E}\kern-.125emX}}
\begin{document}

\title{Prospectivas Tecnológicas de los Sistemas de Información: La Inteligencia Artificial como Núcleo Operativo y Estratégico\\
{\footnotesize Curso: Tendencias de Sistemas de Información | Docente: Cesar Augusto Angulo Calderon}
}

\author{
\IEEEauthorblockN{Mellany K. Apolino Llacta, Victor D. Celadita Romero, John L. Castillo Reupo,}
\IEEEauthorblockN{Jhair R. Melendez Blas y Richard J. Pillaca Machaca}
\IEEEauthorblockA{\textit{Facultad de Ingeniería de Sistemas e Informática} \\
\textit{Universidad Nacional Mayor de San Marcos}\\
Lima, Perú, 2026 \\
\{mellany.apolino, victor.celadita, john.castillo.reupo, jhair.melendez, richard.pillaca\}@unmsm.edu.pe}
}

\maketitle

\begin{abstract}
La Inteligencia Artificial (IA) se ha consolidado como el núcleo funcional de los Sistemas de Información (SI) contemporáneos, transformando infraestructuras reactivas en ecosistemas proactivos. Este documento presenta una revisión sistemática exploratoria sobre las prospectivas tecnológicas de los SI impulsados por IA y la adopción de la Inteligencia de Decisiones (Decision Intelligence). Se analiza la evolución desde modelos orientados a datos (data-driven) hacia sistemas autónomos (AI-driven), evaluando su impacto en la eficiencia operativa, arquitecturas en el borde (Edge Computing) y los retos éticos emergentes. Los hallazgos proporcionan un marco para comprender cómo la automatización de decisiones redefinirá la estrategia organizacional en el mediano y largo plazo.
\end{abstract}

\begin{IEEEkeywords}
Inteligencia Artificial, Sistemas de Información, Decision Intelligence, Machine Learning, Automatización.
\end{IEEEkeywords}

\section{Introducción}
 
La era contemporánea de la disrupción digital ha trascendido la mera digitalización de procesos físicos para situarse en un estado donde la Inteligencia Artificial (IA) actúa como la arquitectura fundamental de los Sistemas de Información (SI), ya que, a lo largo de las últimas décadas, la acumulación masiva de datos, el incremento exponencial de la capacidad computacional y los avances sostenidos en diseño algorítmico han reorientado el campo de la IA casi de forma exclusiva hacia el \textit{Machine Learning} (ML) y los enfoques impulsados por datos \cite{storey2025genai}, consolidando una transición estructural que redefine el rol tecnológico de las organizaciones modernas.
 
En este nuevo paradigma, los SI han dejado de ser repositorios pasivos de información transaccional para convertirse en entidades cognitivas capaces de generar conocimiento autónomo, anticipar comportamientos y prescribir acciones estratégicas~\cite{natta2025,seth2025}, lo que supone una evolución de ecosistemas reactivos a proactivos que impone una reingeniería profunda tanto de las infraestructuras tecnológicas como de la cultura organizacional, dado que los datos ya no son únicamente un recurso de soporte, sino el combustible de un motor automatizado de innovación continua~\cite{murire2024}.
 
Dentro de este contexto, el surgimiento de la \textit{Decision Intelligence} (DI) representa la convergencia entre la predicción algorítmica y la ejecución estratégica accionable, integrando la ciencia de datos con la teoría gerencial y las ciencias sociales para cerrar la brecha existente entre los datos crudos y las decisiones organizacionales de alto impacto~\cite{heilig2024,jeyanthi2022}, cuya adopción no implica la sustitución del juicio humano, sino su \textit{augmentación} sistemática mediante modelos de IA que evalúan múltiples escenarios en tiempo real~\cite{tasleem2024}.
 
El presente trabajo tiene como objetivo proporcionar un análisis
comprensivo de las prospectivas tecnológicas de los Sistemas de
Información centrados en la IA, con base en una revisión sistemática
exploratoria de la literatura científica reciente. Específicamente, se
busca: (i)~caracterizar la transición de los modelos
\textit{data-driven} hacia los ecosistemas \textit{AI-driven}; (ii)~examinar
el rol de la Decision Intelligence como disciplina estratégica emergente;
(iii)~analizar las arquitecturas de Edge Computing como habilitadoras de
la analítica en tiempo real; y (iv)~identificar los desafíos éticos y de
gobernanza que condicionan la escalabilidad de estos sistemas. A través
de este enfoque estructurado, se busca ofrecer una visión del futuro de
los SI que sirva de referencia para investigadores y profesionales que
navegan la era de la Inteligencia Artificial
ubicua~\cite{natta2025,seth2025,tasleem2024}.

\section{Marco Teórico}
 
\subsection{Evolución histórica de los Sistemas de Información}
 
La trayectoria evolutiva de los Sistemas de Información puede
comprenderse como una progresión por fases sucesivas, cada una
caracterizada por cambios cualitativos en la relación entre la tecnología
y la toma de decisiones organizacionales. En su etapa fundacional,
durante las décadas de 1960 y 1970, los SI se concretaban en Sistemas de
Procesamiento de Transacciones (TPS, \textit{Transaction Processing
Systems}) cuya función central era el registro y almacenamiento
automatizado de operaciones rutinarias. Estos sistemas introdujeron la
informatización de procesos administrativos, pero carecían de toda
capacidad analítica~\cite{wu2025aigen}.
 
La década de 1980 marcó la transición hacia los Sistemas de Información
Gerencial (MIS, \textit{Management Information Systems}) y los Sistemas
de Soporte a la Decisión (DSS), que incorporaron rudimentarias capacidades
de síntesis y reporte para apoyar a mandos medios y directivos. Este
período también vio surgir los primeros sistemas expertos, que codificaban
conocimiento especializado mediante reglas lógicas; sin embargo, la
complejidad del mantenimiento de dichas reglas y las limitaciones del
hardware disponible precipitaron el primer «invierno de la IA» en la
segunda mitad de los años
ochenta~\cite{storey2025genai}.
 
El punto de inflexión más significativo ocurre desde finales de la
década de 2000, cuando la confluencia de tres factores —grandes volúmenes
de datos de entrenamiento, la masificación de GPUs de alto rendimiento y
el avance en arquitecturas de redes neuronales profundas— relanza el
campo de la IA bajo el paradigma del \textit{Machine Learning}
(ML)~\cite{janiesch2021ml}. El ML, y particularmente el \textit{Deep
Learning} (DL), habilitó a los SI para procesar volúmenes masivos de
datos heterogéneos e identificar patrones que resultaban invisibles para
el análisis estadístico convencional~\cite{natta2025,seth2025}. Esta
capacidad inauguró la transición desde el análisis descriptivo hacia
modalidades predictivas y prescriptivas, vitales para la competitividad
en mercados volátiles~\cite{janiesch2021ml}.
 
\subsection{Del paradigma \textit{data-driven} al \textit{AI-driven}}
 
La distinción conceptual entre sistemas orientados a datos
(\textit{data-driven}) y sistemas impulsados por IA (\textit{AI-driven})
es fundamental para comprender la prospectiva actual. Mientras que los
primeros utilizan datos históricos como insumo para apoyar decisiones
humanas, los segundos permiten que el propio sistema guíe la trayectoria
organizacional mediante una inteligencia integrada que aprende y se
adapta en tiempo real~\cite{natta2025}. En términos de generaciones
tecnológicas, este salto corresponde al paso de sistemas de IA 1.0
—orientados al procesamiento de información y al reconocimiento de
patrones— hacia arquitecturas de IA 2.0 capaces de planificar,
decidir y actuar con autonomía creciente~\cite{wu2025aigen}.
 
Esta transición no es meramente técnica: implica una reconfiguración
cultural profunda dentro de las organizaciones. La adopción de sistemas
\textit{AI-driven} exige que los profesionales del área operen en la
intersección de la ingeniería de datos, la ciencia de la decisión y la
ética algorítmica~\cite{murire2024}. Las organizaciones que no
logran alinear su cultura interna con este nuevo paradigma enfrentan
resistencia al cambio, brechas de competencias y desafíos de gobernanza
que limitan los beneficios operativos esperados~\cite{murire2024}.
 
\subsection{Componentes tecnológicos centrales}
 
Los sistemas \textit{AI-driven} descansan sobre tres pilares
tecnológicos interdependientes. En primer lugar, las infraestructuras de
datos unificadas —frecuentemente denominadas \textit{Data Lakes}
inteligentes— permiten la ingesta y el procesamiento de datos no
estructurados por parte de redes neuronales profundas, habilitando una
comprensión holística del entorno corporativo~\cite{seth2025}. En segundo
lugar, los algoritmos de ML y DL constituyen el núcleo analítico del
sistema; su capacidad para generalizar a partir de grandes corpus de
datos y detectar relaciones no lineales los posiciona como el estándar
para la interpretación de grandes bases de
datos~\cite{janiesch2021ml,natta2025}. En tercer lugar, las arquitecturas
híbridas de cómputo en la nube y \textit{Edge Computing} permiten que el
procesamiento ocurra de manera distribuida, reduciendo la latencia y
habilitando la toma de decisiones en el punto exacto de interacción con
el cliente~\cite{seth2025}.
 
\subsection{Decision Intelligence como marco teórico integrador}
 
La \textit{Decision Intelligence} (DI) emerge en este contexto como el
marco teórico que orquesta la unión entre la predicción algorítmica y la
acción empresarial~\cite{heilig2024}. A diferencia de las herramientas
de \textit{Business Intelligence} (BI) tradicionales —orientadas a la
generación de reportes descriptivos—, la DI utiliza modelos de IA para
evaluar múltiples escenarios futuros y proponer la ruta de acción con
mayor probabilidad de éxito según los objetivos del
negocio~\cite{heilig2024,jeyanthi2022}. Este enfoque sistémico permite
modelar la incertidumbre y gestionar la complejidad de una manera que las
heurísticas humanas difícilmente podrían igualar~\cite{tasleem2024}.
 
Un hallazgo relevante de la literatura es que el valor diferencial de la
DI no reside únicamente en la optimización algorítmica, sino en su
capacidad para integrar la intuición humana con los modelos predictivos
de IA~\cite{jeyanthi2022}. Este hallazgo anticipa un futuro donde la
toma de decisiones no será un proceso binario (humano \textit{versus}
máquina), sino un ecosistema híbrido donde los sistemas de IA proponen
escenarios optimizados y los líderes organizacionales aportan el contexto
ético, cultural y estratégico que los algoritmos aún no pueden
modelar~\cite{tasleem2024,natta2025}.
 
\subsection{Desafíos éticos y de gobernanza}
 
La autonomía creciente de los sistemas \textit{AI-driven} plantea
interrogantes fundamentales sobre la ética algorítmica y la necesidad de
mantener un control humano significativo. Los sesgos sistémicos, la
opacidad en la toma de decisiones automatizadas y la responsabilidad
distribuida constituyen riesgos que se incrementan en proporción directa
al nivel de autonomía del sistema~\cite{tasleem2024}. La confianza
organizacional en la IA depende, en última instancia, de la capacidad
para auditar, explicar y, cuando sea necesario, revertir las decisiones
automatizadas~\cite{natta2025}.
 
En este sentido, la literatura proyecta el surgimiento de marcos
regulatorios obligatorios para la transparencia algorítmica, especialmente
en sectores críticos como salud, finanzas e infraestructura
pública~\cite{storey2025genai}. La prospectiva apunta hacia sistemas de
\textit{ethical self-monitoring} que no solo detecten y corrijan fallas
técnicas de manera autónoma, sino que también identifiquen y mitiguen
desviaciones éticas en sus propios procesos de decisión~\cite{natta2025}.
Este campo constituye la frontera más avanzada de la investigación en
sistemas de información y tiene implicaciones directas para la formulación
de políticas tecnológicas en el contexto latinoamericano.

\section{Metodología de Búsqueda}
La metodología empleada para este análisis consiste en una revisión sistemática exploratoria fundamentada en las directrices del protocolo PRISMA. Este enfoque garantiza que la selección de la literatura científica sea transparente, reproducible y esté libre de sesgos selectivos \cite{natta2025}. El proceso se inició con la definición de preguntas de investigación críticas que buscan desentrañar el papel de la IA en la arquitectura de los sistemas modernos y su impacto en la toma de decisiones automatizada.

Para la recuperación de los documentos, se consultaron repositorios científicos de prestigio internacional, asegurando una cobertura exhaustiva de los desarrollos más recientes en ingeniería de sistemas y ciencias de la computación. En la Tabla \ref{tab:estrategia} se detallan las bases de datos y la ecuación lógica empleada.

% TABLA 1: ESTRATEGIA DE BÚSQUEDA (ESTILO APA 7)
\begin{table}[htbp]
\caption{Estrategia de Búsqueda y Repositorios}
\label{tab:estrategia}
\centering
\begin{tabular}{p{2.5cm} p{5cm}}
\toprule
\textbf{Elemento} & \textbf{Descripción Aplicada} \\
\midrule
\textit{Bases de Datos} & IEEE Xplore, SpringerLink, Google Académico (Scopus, Semantic Scholar). \\
\addlinespace
\textit{Cadena de} & ("Artificial Intelligence" OR "Machine Learning") AND ("Information Systems" OR "AI-driven") AND ("Decision Intelligence") \\
\textit{Búsqueda} & \\
\bottomrule
\end{tabular}
\end{table}

El proceso de selección de artículos se rigió por criterios estrictos, limitando la búsqueda a publicaciones recientes (2021-2026) para asegurar la vigencia tecnológica de las prospectivas presentadas. La Tabla \ref{tab:criterios} expone los parámetros de filtrado aplicados durante el tamizado de la literatura.

% TABLA 2: CRITERIOS DE INCLUSIÓN Y EXCLUSIÓN (ESTILO APA 7)
\begin{table}[htbp]
\caption{Criterios de Inclusión y Exclusión}
\label{tab:criterios}
\centering
\begin{tabular}{ll p{4.2cm}}
\toprule
\textbf{Categoría} & \textbf{Cód.} & \textbf{Descripción del Criterio} \\
\midrule
\textit{Inclusión} & I01 & Publicaciones en el periodo 2021-2026. \\
 & I02 & Artículos enfocados explícitamente en sistemas AI-driven y Decision Intelligence. \\
 & I03 & Estudios con revisión por pares. \\
\midrule
\textit{Exclusión} & E01 & Literatura gris o blogs corporativos. \\
 & E02 & Artículos puramente algorítmicos sin enfoque organizacional en Sistemas de Información. \\
\bottomrule
\end{tabular}
\end{table}

Tras la identificación inicial de registros mediante la cadena de búsqueda, se realizó un tamizado exhaustivo de títulos y resúmenes para descartar estudios que no situaban a la IA como el núcleo de un sistema de información organizacional. Los documentos seleccionados fueron analizados a texto completo, extrayendo datos clave sobre metodologías, resultados y aportes tecnológicos específicos \cite{natta2025}. Este proceso de filtrado resultó en una muestra final de 5 artículos científicos de alta calidad que sirven como la base de evidencia para el análisis comparativo.

\section{Prospectiva}
La convergencia acelerada entre la Inteligencia Artificial y los Sistemas de Información está delineando escenarios de transformación profunda para las organizaciones en el horizonte 2030. A partir del análisis de las tendencias identificadas en la literatura revisada, es posible proyectar tres escenarios futuros predominantes que redefinirán la arquitectura, la gobernanza y el valor estratégico de los SI.

\subsection{Sistemas de Información Cognitivos y Autónomos}
El primer escenario prospectivo apunta hacia la consolidación de los SI como entidades cognitivas plenamente autónomas. Como señalan Natta (2025) y Seth et al. (2025), la transición de modelos reactivos a ecosistemas proactivos no es una tendencia transitoria, sino una reconfiguración estructural del rol tecnológico en las organizaciones. En este horizonte, los sistemas no solo procesarán datos históricos para generar reportes, sino que establecerán bucles de retroalimentación continua donde la toma de decisiones operativas —y eventualmente estratégicas— será ejecutada por agentes de IA con mínima intervención humana \cite{natta2025, seth2025}.

La materialización de este escenario depende de la maduración de tres pilares tecnológicos: (i) arquitecturas de Edge Computing distribuidas que reduzcan la latencia crítica para decisiones en tiempo real; (ii) modelos de Machine Learning explicables (XAI) que generen confianza en los procesos automatizados; y (iii) infraestructuras de datos unificadas que eliminen los silos organizacionales que hoy fragmentan la inteligencia empresarial \cite{seth2025, tasleem2024}. La investigación de Seth et al. (2025) demuestra que las arquitecturas híbridas edge-cloud ya son técnicamente viables para soportar millones de transacciones con análisis inteligente subyacente, lo que sugiere que la barrera actual no es tecnológica, sino de adopción organizacional y gobernanza de datos.

\subsection{Decision Intelligence como Disciplina Estratégica Central}
El segundo escenario proyecta la institutionalización de la Decision Intelligence (DI) como una disciplina transversal alineada directamente con la alta dirección. Heilig y Scheer (2024) argumentan que la DI representa la evolución natural de la Business Intelligence, cerrando la brecha entre la generación de insights y la ejecución estratégica accionable \cite{heilig2024}. En el mediano plazo, es previsible que las organizaciones establezcan unidades específicas de DI, lideradas por perfiles híbridos que combinen competencias en ciencia de datos, teoría de la decisión y estrategia organizacional.

Jeyanthi et al. (2022) enfatizan que el valor diferencial de la DI no reside únicamente en la optimización algorítmica, sino en su capacidad para integrar la intuición humana con los modelos predictivos de IA \cite{jeyanthi2022}. Este hallazgo sugiere un futuro donde la toma de decisiones no será un proceso binario (humano vs. máquina), sino un ecosistema híbrido donde los sistemas de IA proponen escenarios optimizados y los líderes organizacionales aportan contexto ético, cultural y estratégico que los algoritmos no pueden modelar. La prospectiva, por tanto, no apunta a la sustitución del juicio humano, sino a su augmentación sistemática mediante inteligencia artificial.

\subsection{Gobernanza Ética y Resiliencia Algorítmica}
El tercer escenario —y potencialmente el más disruptivo— concierne a la gobernanza ética y la resiliencia de los sistemas autónomos. Tasleem et al. (2024) advierten que la integración de la intuición humana con modelos de IA introduce vulnerabilidades relacionadas con sesgos algorítmicos, opacidad en la toma de decisiones y responsabilidad distribuida \cite{tasleem2024}. A medida que los SI asuman mayor autonomía, la probabilidad de fallos sistémicos con impacto organizacional o social significativo aumenta exponencialmente.

En respuesta a estos riesgos, la literatura proyecta el surgimiento de marcos regulatorios obligatorios para la transparencia algorítmica, especialmente en sectores críticos como salud, finanzas e infraestructura. Natta (2025) destaca que el éxito de la implementación de DI depende de la capacidad de las organizaciones para establecer mecanismos de gobernanza que equilibren la automatización con la supervisión humana significativa \cite{natta2025}. Este escenario implica que los SI del futuro no serán evaluados únicamente por su eficiencia operativa, sino por su conformidad con estándares éticos y su capacidad de autoauditoría.

La prospectiva de sistemas auto-curativos (\textit{self-healing systems}) mencionada en el marco teórico adquiere aquí una dimensión regulatoria: los SI deberán no solo detectar y corregir fallas técnicas de manera autónoma, sino también identificar y mitigar desviaciones éticas en sus propios procesos de decisión. Esta capacidad de \textit{ethical self-monitoring} representa la frontera más avanzada de la investigación en sistemas de información y constituye un campo de investigación emergente con implicaciones directas para la formulación de políticas tecnológicas.

\subsection{Síntesis de Escenarios}
La interacción dinámica entre estos tres escenarios —autonomía cognitiva, institutionalización de la DI y gobernanza ética— definirá la trayectoria evolutiva de los Sistemas de Información en la próxima década. Las organizaciones que logren orquestar estos tres vectores de manera simultánea obtendrán ventajas competitivas sustentables, mientras que aquellas que descuiden cualquiera de ellos enfrentarán riesgos operativos, reputacionales y regulatorios. La prospectiva tecnológica, en última instancia, no es determinista: su materialización depende de decisiones estratégicas, inversiones en capital humano y la capacidad de las instituciones para adaptar sus estructuras de gobernanza a un entorno donde la inteligencia artificial ya no es una herramienta, sino el núcleo operativo y estratégico del sistema de información.




\section{Conclusiones}
La presente investigación ha permitido establecer un marco comprensivo sobre las prospectivas tecnológicas de los Sistemas de Información en la era de la Inteligencia Artificial. A través de un análisis sistemático de la literatura científica reciente, se han identificado tendencias, arquitecturas y desafíos que configuran el estado actual y futuro de esta disciplina.

En primer lugar, se constata que la transición de sistemas orientados a datos (\textit{data-driven}) hacia ecosistemas impulsados por IA (\textit{AI-driven}) constituye un cambio de paradigma irreversible. Los artículos analizados demuestran de manera convergente que los SI contemporáneos ya no funcionan como repositorios pasivos de información, sino como entidades cognitivas capaces de generar conocimiento autónomo, anticipar comportamientos y prescribir acciones estratégicas \cite{natta2025, seth2025, tasleem2024}. Esta evolución redefine la función del profesional de sistemas de información, quien debe ahora operar en la intersección entre la ingeniería de datos, la ciencia de la decisión y la ética algorítmica.

En segundo lugar, la Decision Intelligence emerge como la disciplina integradora que articula la predicción algorítmica con la ejecución organizacional. Los trabajos de Heilig y Scheer (2024) y Jeyanthi et al. (2022) evidencian que la DI no es una mera extensión tecnológica de la Business Intelligence, sino un marco teórico y práctico que requiere la síntesis de múltiples dominios del conocimiento \cite{heilig2024, jeyanthi2022}. Su implementación efectiva demanda infraestructuras de datos robustas, modelos de IA explicables y, fundamentalmente, una reconfiguración de la cultura organizacional hacia la aceptación de la toma de decisiones augmentada por máquinas.

En tercer lugar, el análisis comparativo revela que las arquitecturas híbridas de Edge Computing y computación en la nube representan el estándar emergente para el despliegue de SI inteligentes. La investigación de Seth et al. (2025) valida que estas arquitecturas no solo resuelven problemas de latencia y escalabilidad, sino que habilitan nuevos modelos de negocio basados en la analítica en tiempo real \cite{seth2025}. Sin embargo, la adopción generalizada de estas arquitecturas enfrenta obstáculos significativos relacionados con la interoperabilidad de sistemas, la seguridad de datos distribuidos y la complejidad de la gobernanza en entornos multi-nube.

En cuarto lugar, los desafíos no resueltos identificados en la literatura —particularmente los relacionados con sesgos algorítmicos, transparencia y responsabilidad en sistemas autónomos— constituyen barreras críticas para la materialización plena de las prospectivas tecnológicas descritas. Tasleem et al. (2024) y Natta (2025) coinciden en señalar que la confianza en los sistemas de IA depende de la capacidad para auditar, explicar y, cuando sea necesario, revertir decisiones automatizadas \cite{tasleem2024, natta2025}. La ausencia de marcos regulatorios uniformes y de estándares técnicos de explicabilidad representa, en la actualidad, el cuello de botella más significativo para la escalabilidad ética de estos sistemas.

Finalmente, la prospectiva desarrollada en este trabajo proyecta un horizonte donde los Sistemas de Información alcanzarán niveles de autonomía cognitiva y auto-gestión sin precedentes. No obstante, esta autonomía tecnológica no implica la irrelevancia del factor humano. Por el contrario, la evidencia analizada sugiere que el valor competitivo futuro residirá en la capacidad de las organizaciones para diseñar arquitecturas de decisión híbridas que potencien las fortalezas complementarias de la inteligencia humana y la artificial. En este sentido, los Sistemas de Información del futuro serán tan inteligentes como lo permitan sus mecanismos de gobernanza ética y su diseño centrado en la augmentación, más que en la sustitución, de la capacidad humana de juicio.

Este trabajo contribuye al campo de las tendencias de sistemas de información al proporcionar una visión integrada que vincula avances tecnológicos con implicaciones organizacionales y éticas. Se recomienda que investigaciones futuras profundicen en el desarrollo de métricas de madurez para la adopción de Decision Intelligence, en el diseño de frameworks de gobernanza algorítmica aplicables a contextos latinoamericanos y en la evaluación empírica del impacto de sistemas AI-driven en indicadores de desempeño organizacional de mediano y largo plazo.

% BIBLIOGRAFÍA CON FORMATO DE TEXTO APA 7 (Pero manteniendo la numeración IEEE)
\begin{thebibliography}{00}

\bibitem{heilig2024} Heilig, T., \& Scheer, I. (2024). Decision Intelligence: Making Relevant Information Visible and Actionable. \textit{IEEE Xplore}. https://doi.org/10.1109/10952523

\bibitem{jeyanthi2022} Jeyanthi, P. M., Hack-Polay, D., \& Choudhury, T. (2022). The Rise of Decision Intelligence: AI That Optimizes Decision-Making. En \textit{Decision Intelligence Analytics and the Implementation of Strategic Business Management} (pp. 55--68). Springer. https://doi.org/10.1007/978-3-030-82763-2\_4

\bibitem{natta2025} Natta, P. K. (2025). AI-Driven Decision Intelligence: Optimizing Enterprise Strategy with AI-Augmented Insights. \textit{Journal of Computer Science and Technology Studies, 7}(2), 146--152. https://doi.org/10.32996/jcsts.2025.7.2.13

\bibitem{tasleem2024} Tasleem, N., Raghav, R. S., Ansari, M. N., \& Sharma, A. J. (2024). A Decision Intelligence Framework: Integrating Human Intuition with AI Models. \textit{Journal of Artificial Intelligence General Science, 7}(1), 304--319. 

\bibitem{seth2025} Seth, D. K., Ratra, K. K., \& Sundareswaran, A. P. (2025). AI driven hybrid edge cloud architecture for real time big data analytics and scalable communication. \textit{IEEE SoutheastCon 2025}. 

\bibitem{murire2024} Murire, O.~T. (2024). Artificial Intelligence and Its Role in Shaping Organizational Work Practices and Culture. \textit{Administrative Sciences}, \textit{14}(12), 316. https://doi.org/10.3390/admsci14120316

\bibitem{storey2025genai} Storey, V.~C., Yue, W.~T., Zhao, J.~L., \& Lukyanenko, R. (2025). Generative Artificial Intelligence: Evolving Technology, Growing Societal Impact, and Opportunities for Information Systems Research. \textit{Information Systems Frontiers}. https://doi.org/10.1007/s10796-025-10581-7

\bibitem{wu2025aigen} Wu, J., You, H., \& Du, J. (2025). AI Generations: From AI 1.0 to AI 4.0. \textit{Frontiers in Artificial Intelligence}, \textit{8}, 1585629. https://doi.org/10.3389/frai.2025.1585629

\bibitem{janiesch2021ml} Janiesch, C., Zschech, P., \& Heinrich, K. (2021). Machine learning and deep learning. \textit{Electronic Markets}, \textit{31}, 685--695. https://doi.org/10.1007/s12525-021-00475-2

\end{thebibliography}

\end{document}